% ======================================================================
%  UPPGIFT 1: SPÄNNANDE TRÄD
% ======================================================================
\section{1 Spännande träd}

\uppgiftsbild{1}{Spännande träd}
{köra Prim och Kruskal för hand — och se \emph{varför} två olika giriga val båda blir rätt.}
{\begin{itemize}\setlength{\itemsep}{1pt}
  \item Kandidatmängden: \emph{hela grafen} (Kruskal) eller \emph{kanten ut ur trädet} (Prim).
  \item Tillståndet under körningen: skog eller träd.
  \item Vad MST beror på: bara vikternas \emph{ordning} (kontrast: kortaste väg).
\end{itemize}}

\begin{frame}{Uppgift 1: Spännande träd}
\begin{columns}[T]
\begin{column}{0.44\textwidth}
\centering
\begin{tikzpicture}[scale=0.78]
  \node[gv] (a) at (3.1,5.2) {$a$};
  \node[gv] (b) at (2.1,4.2) {$b$};
  \node[gv] (c) at (4.3,4.1) {$c$};
  \node[gv] (d) at (1.0,3.1) {$d$};
  \node[gv] (e) at (3.1,3.1) {$e$};
  \node[gv] (f) at (4.7,2.1) {$f$};
  \node[gv] (g) at (1.0,1.0) {$g$};
  \node[gv] (h) at (3.1,1.0) {$h$};
  \node[gv] (i) at (0.0,0.0) {$i$};
  \begin{scope}[on background layer]
  \draw[ge] (a) -- node[gw,above right]{1} (c);
  \draw[ge] (b) -- node[gw,above right]{2} (e);
  \draw[ge] (e) -- node[gw,left]{3} (h);
  \draw[ge] (g) -- node[gw,below]{4} (h);
  \draw[ge] (b) -- node[gw,above left]{5} (d);
  \draw[ge] (d) -- node[gw,right]{6} (g);
  \draw[ge] (a) .. controls (-0.2,5.4) and (-0.4,1.3) .. node[gw,left,pos=0.5]{7} (g);
  \draw[ge] (g) -- node[gw,above left]{8} (i);
  \draw[ge] (a) -- node[gw,above left]{9} (b);
  \draw[ge] (e) -- node[gw,above right]{10} (f);
  \end{scope}
\end{tikzpicture}

\end{column}
\begin{column}{0.53\textwidth}
\begin{infobox}
Visa hur ett minimalt spännande träd hittas i grafen med både \textbf{Prims} och
\textbf{Kruskals} algoritmer.
\end{infobox}
\small
\begin{itemize}\setlength{\itemsep}{3pt}
  \item \textbf{Spännande träd:} delgraf som är ett träd och når \emph{alla} hörn.
        Alltså exakt $|V|-1$ kanter; här $9-1=8$.
  \item \textbf{Minimalt:} minsta möjliga summa av kantvikterna.
  \item Båda algoritmerna är \emph{giriga}. De är bara giriga på olika saker.
\end{itemize}
\end{column}
\end{columns}
\end{frame}

\begin{frame}{Två giriga algoritmer — och en minnesregel}
\begin{columns}[T]
\begin{column}{0.49\textwidth}
\begin{infobox}[title={\textbf{Kruskal}},coltitle=white,colbacktitle=kthblue,fonttitle=\small,
  toptitle=1pt,bottomtitle=1pt]
Ta den \textbf{globalt} billigaste kanten, i hela grafen, som förbinder två
\emph{olika komponenter}. Tillståndet är en \emph{skog}.
\end{infobox}
\end{column}
\begin{column}{0.49\textwidth}
\begin{infobox}[title={\textbf{Prim}},coltitle=white,colbacktitle=kthblue,fonttitle=\small,
  toptitle=1pt,bottomtitle=1pt]
Ta den billigaste kanten som går \textbf{ut ur trädet}. Tillståndet är ett enda
\emph{träd} som växer från en rot.
\end{infobox}
\end{column}
\end{columns}
\medskip
\begin{center}
\begin{tikzpicture}
  \node[draw=kthorange,line width=1pt,fill=lampyellow!20,rounded corners=4pt,inner sep=7pt,
        font=\LARGE] (m) {Krusk\textcolor{kthorange}{\textbf{AL}}\ \ är\ \ Glob\textcolor{kthorange}{\textbf{AL}}};
  \node[font=\footnotesize,text=black!60,below=3pt of m]
     {Prim är lokal. (Rimmar inte. Det är därför ingen kommer ihåg det.)};
\end{tikzpicture}
\end{center}
\varstrip{Kontrast}{grafen, girigheten}{kandidatmängden}{global kantlista mot lokal snittkant}
\end{frame}

\begin{frame}{Pseudokod}
\begin{columns}[T]
\begin{column}{0.48\textwidth}
{\small\bfseries\color{kthnavy}Kruskal$(G,w)$}
\footnotesize
\begin{algorithmic}[1]
\State $T\gets\emptyset$
\ForAll{$v\in V$} \State \textsc{MakeSet}$(v)$ \EndFor
\State sortera $E$ i växande vikt
\ForAll{$(u,v)\in E$ i sorterad ordning}
  \If{$\textsc{Find}(u)\neq\textsc{Find}(v)$}
    \State $T\gets T\cup\{(u,v)\}$
    \State \textsc{Union}$(u,v)$
  \EndIf
\EndFor
\State \Return $T$
\end{algorithmic}
\end{column}
\begin{column}{0.5\textwidth}
{\small\bfseries\color{kthnavy}Prim$(G,w,r)$}
\footnotesize
\begin{algorithmic}[1]
\ForAll{$v\in V$} \State $\mathit{key}[v]\gets\infty$;\ $\pi[v]\gets\textsc{nil}$ \EndFor
\State $\mathit{key}[r]\gets0$
\State $Q\gets V$ \Comment{prioritet $\mathit{key}$}
\While{$Q\neq\emptyset$}
  \State $u\gets\textsc{ExtractMin}(Q)$
  \ForAll{$(u,v)\in E$}
    \If{\colorbox{lampyellow!50}{$v\in Q$} \textbf{och} $w(u,v)<\mathit{key}[v]$}
      \State $\pi[v]\gets u$;\ $\mathit{key}[v]\gets w(u,v)$ 
    \EndIf
  \EndFor
\EndWhile
\State \Return $\{(\pi[v],v) : v\neq r\}$
\end{algorithmic}
\end{column}
\end{columns}
\vfill
\begin{missbox}
”Villkoret $v\in Q$ är en detalj.” Utan det kan ett redan utplockat hörn få ny nyckel
och ny förälder — och då är resultatet inte längre ett träd. (Just den buggen finns i
äldre övningsbilder.)
\end{missbox}
\end{frame}

\begin{frame}{Kruskal på uppgiftens graf}
\begin{columns}[T]
\begin{column}{0.47\textwidth}
\centering
\begin{tikzpicture}[scale=0.76]
  \node[gv,onslide=<2-7>{comp1},onslide=<8->{comp2}] (a) at (3.1,5.2) {$a$};
  \node[gv,onslide=<3->{comp2}] (b) at (2.1,4.2) {$b$};
  \node[gv,onslide=<2-7>{comp1},onslide=<8->{comp2}] (c) at (4.3,4.1) {$c$};
  \node[gv,onslide=<6->{comp2}] (d) at (1.0,3.1) {$d$};
  \node[gv,onslide=<3->{comp2}] (e) at (3.1,3.1) {$e$};
  \node[gv,onslide=<11->{comp2}] (f) at (4.7,2.1) {$f$};
  \node[gv,onslide=<5->{comp2}] (g) at (1.0,1.0) {$g$};
  \node[gv,onslide=<4->{comp2}] (h) at (3.1,1.0) {$h$};
  \node[gv,onslide=<9->{comp2}] (i) at (0.0,0.0) {$i$};
  \begin{scope}[on background layer]
  \draw[kthorange!55,line width=7pt,line cap=round,visible on=<7>] (h) -- (g);
  \draw[kthorange!55,line width=7pt,line cap=round,visible on=<7>] (e) -- (h);
  \draw[kthorange!55,line width=7pt,line cap=round,visible on=<7>] (b) -- (e);
  \draw[kthorange!55,line width=7pt,line cap=round,visible on=<7>] (d) -- (b);
  \draw[kthorange!55,line width=7pt,line cap=round,visible on=<10>] (e) -- (b);
  \draw[kthorange!55,line width=7pt,line cap=round,visible on=<10>] (h) -- (e);
  \draw[kthorange!55,line width=7pt,line cap=round,visible on=<10>] (g) -- (h);
  \draw[kthorange!55,line width=7pt,line cap=round,visible on=<10>] (a) .. controls (-0.2,5.4) and (-0.4,1.3) .. (g);
  \draw[ge,onslide=<2>{nyss},onslide=<3->{mst}] (a) -- node[gw,above right]{1} (c);
  \draw[ge,onslide=<3>{nyss},onslide=<4->{mst}] (b) -- node[gw,above right]{2} (e);
  \draw[ge,onslide=<4>{nyss},onslide=<5->{mst}] (e) -- node[gw,left]{3} (h);
  \draw[ge,onslide=<5>{nyss},onslide=<6->{mst}] (g) -- node[gw,below]{4} (h);
  \draw[ge,onslide=<6>{nyss},onslide=<7->{mst}] (b) -- node[gw,above left]{5} (d);
  \draw[ge,onslide=<7>{rejnu},onslide=<8->{rej}] (d) -- node[gw,right]{6} (g);
  \draw[ge,onslide=<8>{nyss},onslide=<9->{mst}] (a) .. controls (-0.2,5.4) and (-0.4,1.3) .. node[gw,left,pos=0.5]{7} (g);
  \draw[ge,onslide=<9>{nyss},onslide=<10->{mst}] (g) -- node[gw,above left]{8} (i);
  \draw[ge,onslide=<10>{rejnu},onslide=<11->{rej}] (a) -- node[gw,above left]{9} (b);
  \draw[ge,onslide=<11->{nyss}] (e) -- node[gw,above right]{10} (f);
  \end{scope}
\end{tikzpicture}
\end{column}
\begin{column}{0.5\textwidth}
\begin{tikzpicture}[x=1cm,y=-0.40cm,font=\small]
  \node[anchor=west,font=\footnotesize\bfseries,text=kthnavy] at (-0.35,-0.9) {Alla kanter, globalt sorterade};
  \fill[kthorange!22,visible on=<2>] (-0.35,0-0.45) rectangle (5.6,0+0.45);
  \node[visible on=<2>,text=kthorange] at (-0.18,0) {$\blacktriangleright$};
  \node[anchor=east] at (0.55,0) {1};
  \node[anchor=west] at (0.75,0) {$a$--$c$};
  \node[anchor=west,visible on=<2->] at (1.75,0) {\ja};
  \fill[kthorange!22,visible on=<3>] (-0.35,1-0.45) rectangle (5.6,1+0.45);
  \node[visible on=<3>,text=kthorange] at (-0.18,1) {$\blacktriangleright$};
  \node[anchor=east] at (0.55,1) {2};
  \node[anchor=west] at (0.75,1) {$b$--$e$};
  \node[anchor=west,visible on=<3->] at (1.75,1) {\ja};
  \fill[kthorange!22,visible on=<4>] (-0.35,2-0.45) rectangle (5.6,2+0.45);
  \node[visible on=<4>,text=kthorange] at (-0.18,2) {$\blacktriangleright$};
  \node[anchor=east] at (0.55,2) {3};
  \node[anchor=west] at (0.75,2) {$e$--$h$};
  \node[anchor=west,visible on=<4->] at (1.75,2) {\ja};
  \fill[kthorange!22,visible on=<5>] (-0.35,3-0.45) rectangle (5.6,3+0.45);
  \node[visible on=<5>,text=kthorange] at (-0.18,3) {$\blacktriangleright$};
  \node[anchor=east] at (0.55,3) {4};
  \node[anchor=west] at (0.75,3) {$g$--$h$};
  \node[anchor=west,visible on=<5->] at (1.75,3) {\ja};
  \fill[kthorange!22,visible on=<6>] (-0.35,4-0.45) rectangle (5.6,4+0.45);
  \node[visible on=<6>,text=kthorange] at (-0.18,4) {$\blacktriangleright$};
  \node[anchor=east] at (0.55,4) {5};
  \node[anchor=west] at (0.75,4) {$b$--$d$};
  \node[anchor=west,visible on=<6->] at (1.75,4) {\ja};
  \fill[kthorange!22,visible on=<7>] (-0.35,5-0.45) rectangle (5.6,5+0.45);
  \node[visible on=<7>,text=kthorange] at (-0.18,5) {$\blacktriangleright$};
  \node[anchor=east] at (0.55,5) {6};
  \node[anchor=west] at (0.75,5) {$d$--$g$};
  \node[anchor=west,visible on=<7->,font=\scriptsize,text=brick] at (1.75,5) {\nej\ samma komponent $\Rightarrow$ cykel};
  \fill[kthorange!22,visible on=<8>] (-0.35,6-0.45) rectangle (5.6,6+0.45);
  \node[visible on=<8>,text=kthorange] at (-0.18,6) {$\blacktriangleright$};
  \node[anchor=east] at (0.55,6) {7};
  \node[anchor=west] at (0.75,6) {$a$--$g$};
  \node[anchor=west,visible on=<8->] at (1.75,6) {\ja};
  \fill[kthorange!22,visible on=<9>] (-0.35,7-0.45) rectangle (5.6,7+0.45);
  \node[visible on=<9>,text=kthorange] at (-0.18,7) {$\blacktriangleright$};
  \node[anchor=east] at (0.55,7) {8};
  \node[anchor=west] at (0.75,7) {$g$--$i$};
  \node[anchor=west,visible on=<9->] at (1.75,7) {\ja};
  \fill[kthorange!22,visible on=<10>] (-0.35,8-0.45) rectangle (5.6,8+0.45);
  \node[visible on=<10>,text=kthorange] at (-0.18,8) {$\blacktriangleright$};
  \node[anchor=east] at (0.55,8) {9};
  \node[anchor=west] at (0.75,8) {$a$--$b$};
  \node[anchor=west,visible on=<10->,font=\scriptsize,text=brick] at (1.75,8) {\nej\ samma komponent $\Rightarrow$ cykel};
  \fill[kthorange!22,visible on=<11>] (-0.35,9-0.45) rectangle (5.6,9+0.45);
  \node[visible on=<11>,text=kthorange] at (-0.18,9) {$\blacktriangleright$};
  \node[anchor=east] at (0.55,9) {10};
  \node[anchor=west] at (0.75,9) {$e$--$f$};
  \node[anchor=west,visible on=<11->] at (1.75,9) {\ja};
\end{tikzpicture}

\vskip2pt
{\footnotesize\color{kthblue}\only<1>{Komponenter: 9\quad Kanter i $T$: 0\quad Vikt: 0}\only<2>{Komponenter: 8\quad Kanter i $T$: 1\quad Vikt: 1}\only<3>{Komponenter: 7\quad Kanter i $T$: 2\quad Vikt: 3}\only<4>{Komponenter: 6\quad Kanter i $T$: 3\quad Vikt: 6}\only<5>{Komponenter: 5\quad Kanter i $T$: 4\quad Vikt: 10}\only<6>{Komponenter: 4\quad Kanter i $T$: 5\quad Vikt: 15}\only<7>{Komponenter: 4\quad Kanter i $T$: 5\quad Vikt: 15}\only<8>{Komponenter: 3\quad Kanter i $T$: 6\quad Vikt: 22}\only<9>{Komponenter: 2\quad Kanter i $T$: 7\quad Vikt: 30}\only<10>{Komponenter: 2\quad Kanter i $T$: 7\quad Vikt: 30}\only<11>{Komponenter: 1\quad Kanter i $T$: 8\quad Vikt: 40}}
\end{column}
\end{columns}
\vfill
\begin{tcolorbox}[enhanced,colback=kthlight!50,colframe=kthsky,boxrule=0.4pt,arc=2pt,
  left=4pt,right=4pt,top=1pt,bottom=1pt,fontupper=\small,height=0.95cm,valign=center]
\only<1>{Sortera \emph{alla} kanter. Varje hörn är en egen komponent.}\only<2>{$a$--$c$ förbinder två olika komponenter $\Rightarrow$ ta den.}\only<3>{$b$--$e$: på andra sidan grafen. Kruskal bryr sig inte om \emph{var} kanten är — bara att den är billigast \emph{globalt}.}\only<4>{$e$--$h$ $\Rightarrow$ ta den.}\only<5>{$g$--$h$ $\Rightarrow$ ta den.}\only<6>{$b$--$d$ $\Rightarrow$ ta den.}\only<7>{$d$--$g$: $d$ och $g$ hänger redan ihop via $d$--$b$--$e$--$h$--$g$. Kanten skulle sluta en cykel.}\only<8>{$a$--$g$ förenar $\{a,c\}$ med den stora komponenten.}\only<9>{$g$--$i$ $\Rightarrow$ ta den.}\only<10>{$a$--$b$: redan förbundna via $a$--$g$--$h$--$e$--$b$ $\Rightarrow$ cykel.}\only<11>{$e$--$f$ $\Rightarrow$ ta den. $|V|-1=8$ kanter, vikt $40$. Klart.}
\end{tcolorbox}
\end{frame}


\begin{frame}{Samma tillstånd, olika kandidater}
\begin{columns}[T]
\begin{column}{0.49\textwidth}
\centering
{\small\bfseries\color{kthblue}Kruskal efter första kanten}\par\smallskip
\begin{tikzpicture}[scale=0.56,every node/.append style={font=\footnotesize}]
  \node[gv,comp1] (a) at (3.1,5.2) {$a$};
  \node[gv] (b) at (2.1,4.2) {$b$};
  \node[gv,comp1] (c) at (4.3,4.1) {$c$};
  \node[gv] (d) at (1.0,3.1) {$d$};
  \node[gv] (e) at (3.1,3.1) {$e$};
  \node[gv] (f) at (4.7,2.1) {$f$};
  \node[gv] (g) at (1.0,1.0) {$g$};
  \node[gv] (h) at (3.1,1.0) {$h$};
  \node[gv] (i) at (0.0,0.0) {$i$};
  \begin{scope}[on background layer]
  \draw[mst] (a) -- node[gw,above right]{1} (c);
  \draw[nyss] (b) -- node[gw,above right]{2} (e);
  \draw[cand] (e) -- node[gw,left]{3} (h);
  \draw[cand] (g) -- node[gw,below]{4} (h);
  \draw[cand] (b) -- node[gw,above left]{5} (d);
  \draw[cand] (d) -- node[gw,right]{6} (g);
  \draw[cand] (a) .. controls (-0.2,5.4) and (-0.4,1.3) .. node[gw,left,pos=0.5]{7} (g);
  \draw[cand] (g) -- node[gw,above left]{8} (i);
  \draw[cand] (a) -- node[gw,above left]{9} (b);
  \draw[cand] (e) -- node[gw,above right]{10} (f);
  \end{scope}
\end{tikzpicture}
\par
\footnotesize Kandidater: \emph{alla} kanter mellan olika komponenter — här alla som återstår.\\
Väljer \textbf{2} ($b$--$e$) — långt bort. Krusk\textbf{AL}: glob\textbf{AL}.
\end{column}
\begin{column}{0.49\textwidth}
\centering
{\small\bfseries\color{kthblue}Prim efter första kanten}\par\smallskip
\input{gen/globlok-prim}\par
\footnotesize Kandidater: bara kanter \emph{ut ur trädet} $\{a,c\}$: 7 och 9.\\
Väljer \textbf{7} ($a$--$g$), trots att 2 finns. Prim: lokal.
\end{column}
\end{columns}
\varstrip{Kontrast}{grafen och första steget}{algoritmen}{vilken mängd ”billigast” gäller}
\end{frame}

\begin{frame}{Prim på uppgiftens graf (rot $a$)}
\begin{columns}[T]
\begin{column}{0.47\textwidth}
\centering
\begin{tikzpicture}[scale=0.76]
  \node[gv,onslide=<1>{draw=kthorange,line width=1.1pt},onslide=<2->{intree}] (a) at (3.1,5.2) {$a$};
  \node[gv,onslide=<2-6>{draw=kthorange,line width=1.1pt},onslide=<7->{intree}] (b) at (2.1,4.2) {$b$};
  \node[gv,onslide=<2>{draw=kthorange,line width=1.1pt},onslide=<3->{intree}] (c) at (4.3,4.1) {$c$};
  \node[gv,onslide=<4-7>{draw=kthorange,line width=1.1pt},onslide=<8->{intree}] (d) at (1.0,3.1) {$d$};
  \node[gv,onslide=<5>{draw=kthorange,line width=1.1pt},onslide=<6->{intree}] (e) at (3.1,3.1) {$e$};
  \node[gv,onslide=<6-9>{draw=kthorange,line width=1.1pt},onslide=<10->{intree}] (f) at (4.7,2.1) {$f$};
  \node[gv,onslide=<2-3>{draw=kthorange,line width=1.1pt},onslide=<4->{intree}] (g) at (1.0,1.0) {$g$};
  \node[gv,onslide=<4>{draw=kthorange,line width=1.1pt},onslide=<5->{intree}] (h) at (3.1,1.0) {$h$};
  \node[gv,onslide=<4-8>{draw=kthorange,line width=1.1pt},onslide=<9->{intree}] (i) at (0.0,0.0) {$i$};
  \begin{scope}[on background layer]
  \draw[ge,onslide=<2>{cand},onslide=<3>{nyss},onslide=<4->{mst}] (a) -- node[gw,above right]{1} (c);
  \draw[ge,onslide=<6>{cand},onslide=<7>{nyss},onslide=<8->{mst}] (b) -- node[gw,above right]{2} (e);
  \draw[ge,onslide=<5>{cand},onslide=<6>{nyss},onslide=<7->{mst}] (e) -- node[gw,left]{3} (h);
  \draw[ge,onslide=<4>{cand},onslide=<5>{nyss},onslide=<6->{mst}] (g) -- node[gw,below]{4} (h);
  \draw[ge,onslide=<7>{cand},onslide=<8>{nyss},onslide=<9->{mst}] (b) -- node[gw,above left]{5} (d);
  \draw[ge,onslide=<4-6>{cand},onslide=<7->{faint}] (d) -- node[gw,right]{6} (g);
  \draw[ge,onslide=<2-3>{cand},onslide=<4>{nyss},onslide=<5->{mst}] (a) .. controls (-0.2,5.4) and (-0.4,1.3) .. node[gw,left,pos=0.5]{7} (g);
  \draw[ge,onslide=<4-8>{cand},onslide=<9>{nyss},onslide=<10->{mst}] (g) -- node[gw,above left]{8} (i);
  \draw[ge,onslide=<2-5>{cand},onslide=<6->{faint}] (a) -- node[gw,above left]{9} (b);
  \draw[ge,onslide=<6-9>{cand},onslide=<10->{nyss}] (e) -- node[gw,above right]{10} (f);
  \end{scope}
\end{tikzpicture}
\end{column}
\begin{column}{0.5\textwidth}
{\footnotesize\bfseries\color{kthnavy}Prioritetskö $Q$ (nyckel = billigaste kant in i trädet)}\par\smallskip
\footnotesize
\only<1>{%
\begin{tabular}{@{}l@{\quad}r@{\quad}l@{}}
\toprule hörn & key & $\pi$\\\midrule
$a$ & $0$ & --\\
\multicolumn{3}{@{}l}{\color{black!50}+ 8 hörn med nyckel $\infty$}\\
\bottomrule
\end{tabular}\par\smallskip
{\scriptsize\color{kthblue}Utplockade: --}}
\only<2>{%
\begin{tabular}{@{}l@{\quad}r@{\quad}l@{}}
\toprule hörn & key & $\pi$\\\midrule
$c$ & $1$ & $a$\ \textcolor{kthgrass}{(ny)}\\
$g$ & $7$ & $a$\ \textcolor{kthgrass}{(ny)}\\
$b$ & $9$ & $a$\ \textcolor{kthgrass}{(ny)}\\
\multicolumn{3}{@{}l}{\color{black!50}+ 5 hörn med nyckel $\infty$}\\
\bottomrule
\end{tabular}\par\smallskip
{\scriptsize\color{kthblue}Utplockade: $a$}}
\only<3>{%
\begin{tabular}{@{}l@{\quad}r@{\quad}l@{}}
\toprule hörn & key & $\pi$\\\midrule
$g$ & $7$ & $a$\\
$b$ & $9$ & $a$\\
\multicolumn{3}{@{}l}{\color{black!50}+ 5 hörn med nyckel $\infty$}\\
\bottomrule
\end{tabular}\par\smallskip
{\scriptsize\color{kthblue}Utplockade: $a$, $c$}}
\only<4>{%
\begin{tabular}{@{}l@{\quad}r@{\quad}l@{}}
\toprule hörn & key & $\pi$\\\midrule
$h$ & $4$ & $g$\ \textcolor{kthgrass}{(ny)}\\
$d$ & $6$ & $g$\ \textcolor{kthgrass}{(ny)}\\
$i$ & $8$ & $g$\ \textcolor{kthgrass}{(ny)}\\
$b$ & $9$ & $a$\\
\multicolumn{3}{@{}l}{\color{black!50}+ 2 hörn med nyckel $\infty$}\\
\bottomrule
\end{tabular}\par\smallskip
{\scriptsize\color{kthblue}Utplockade: $a$, $c$, $g$}}
\only<5>{%
\begin{tabular}{@{}l@{\quad}r@{\quad}l@{}}
\toprule hörn & key & $\pi$\\\midrule
$e$ & $3$ & $h$\ \textcolor{kthgrass}{(ny)}\\
$d$ & $6$ & $g$\\
$i$ & $8$ & $g$\\
$b$ & $9$ & $a$\\
\multicolumn{3}{@{}l}{\color{black!50}+ 1 hörn med nyckel $\infty$}\\
\bottomrule
\end{tabular}\par\smallskip
{\scriptsize\color{kthblue}Utplockade: $a$, $c$, $g$, $h$}}
\only<6>{%
\begin{tabular}{@{}l@{\quad}r@{\quad}l@{}}
\toprule hörn & key & $\pi$\\\midrule
$b$ & $2$ & $e$\ \textcolor{brick}{(var 9)}\\
$d$ & $6$ & $g$\\
$i$ & $8$ & $g$\\
$f$ & $10$ & $e$\ \textcolor{kthgrass}{(ny)}\\

\bottomrule
\end{tabular}\par\smallskip
{\scriptsize\color{kthblue}Utplockade: $a$, $c$, $g$, $h$, $e$}}
\only<7>{%
\begin{tabular}{@{}l@{\quad}r@{\quad}l@{}}
\toprule hörn & key & $\pi$\\\midrule
$d$ & $5$ & $b$\ \textcolor{brick}{(var 6)}\\
$i$ & $8$ & $g$\\
$f$ & $10$ & $e$\\

\bottomrule
\end{tabular}\par\smallskip
{\scriptsize\color{kthblue}Utplockade: $a$, $c$, $g$, $h$, $e$, $b$}}
\only<8>{%
\begin{tabular}{@{}l@{\quad}r@{\quad}l@{}}
\toprule hörn & key & $\pi$\\\midrule
$i$ & $8$ & $g$\\
$f$ & $10$ & $e$\\

\bottomrule
\end{tabular}\par\smallskip
{\scriptsize\color{kthblue}Utplockade: $a$, $c$, $g$, $h$, $e$, $b$, $d$}}
\only<9>{%
\begin{tabular}{@{}l@{\quad}r@{\quad}l@{}}
\toprule hörn & key & $\pi$\\\midrule
$f$ & $10$ & $e$\\

\bottomrule
\end{tabular}\par\smallskip
{\scriptsize\color{kthblue}Utplockade: $a$, $c$, $g$, $h$, $e$, $b$, $d$, $i$}}
\only<10>{%
\begin{tabular}{@{}l@{\quad}r@{\quad}l@{}}
\toprule hörn & key & $\pi$\\\midrule
\multicolumn{3}{@{}l}{\color{black!50}(tom)}\\

\bottomrule
\end{tabular}\par\smallskip
{\scriptsize\color{kthblue}Utplockade: $a$, $c$, $g$, $h$, $e$, $b$, $d$, $i$, $f$}}
\end{column}
\end{columns}
\vfill
\begin{tcolorbox}[enhanced,colback=kthlight!50,colframe=kthsky,boxrule=0.4pt,arc=2pt,
  left=4pt,right=4pt,top=1pt,bottom=1pt,fontupper=\small,height=0.95cm,valign=center]
\only<1>{Starta i $a$: $\mathit{key}[a]=0$, alla andra $\infty$. Alla hörn ligger i kön.}\only<2>{Plocka ut $a$. Grannarna får nycklar: $c$ (1), $g$ (7), $b$ (9).}\only<3>{Plocka ut $c$ (1). Inga nya grannar. Kön är lokal: bara kanter \emph{ut ur trädet} räknas.}\only<4>{Plocka ut $g$ (7). Nya nycklar: $h$ (4), $d$ (6), $i$ (8).}\only<5>{Plocka ut $h$ (4). Ny nyckel: $e$ (3).}\only<6>{Plocka ut $e$ (3). \textsc{Decrease-Key}: $b$ går från 9 till 2. Ny: $f$ (10).}\only<7>{Plocka ut $b$ (2). \textsc{Decrease-Key}: $d$ går från 6 till 5.}\only<8>{Plocka ut $d$ (5). Kanten $d$--$g$ (6) blev aldrig trädkant.}\only<9>{Plocka ut $i$ (8).}\only<10>{Plocka ut $f$ (10). Vikt $1+7+4+3+2+5+8+10=40$. Samma träd som Kruskal (unika vikter!).}
\end{tcolorbox}
\end{frame}


\begin{frame}{Facit och kontroll}
\begin{columns}[T]
\begin{column}{0.44\textwidth}
\centering
\input{gen/graf-mst}
\end{column}
\begin{column}{0.53\textwidth}
\small
\begin{itemize}\setlength{\itemsep}{3pt}
  \item Kanter: $1,2,3,4,5,7,8,10$. \quad Vikt: $\mathbf{40}$.
  \item Kruskal förkastar \textbf{6} och \textbf{9}: båda sluter cykler.
  \item Prim från $a$ plockar i ordningen $a,c,g,h,e,b,d,i,f$.
  \item Samma träd, olika \emph{ordning}. Det är ingen slump: med unika vikter
        finns bara ett MST.
\end{itemize}
\begin{ahabox}
\textbf{Kontroll som alltid fungerar:} rätt antal kanter ($|V|-1$), ingen cykel,
och varje förkastad kant är den tyngsta på sin cykel.
\end{ahabox}
\end{column}
\end{columns}
\end{frame}

% ---------------------------------------------------------------- kaninhål
{\kaninbg
\begin{frame}{\kh Varför fungerar girigheten? Snitt och cykler}
\begin{columns}[T]
\begin{column}{0.49\textwidth}
\centering
{\small\bfseries\color{plum}Snittegenskapen}\par
\begin{tikzpicture}[scale=0.56,every node/.append style={font=\footnotesize}]
  \node[gv] (a) at (3.1,5.2) {$a$};
  \node[gv] (b) at (2.1,4.2) {$b$};
  \node[gv] (c) at (4.3,4.1) {$c$};
  \node[gv] (d) at (1.0,3.1) {$d$};
  \node[gv] (e) at (3.1,3.1) {$e$};
  \node[gv] (f) at (4.7,2.1) {$f$};
  \node[gv] (g) at (1.0,1.0) {$g$};
  \node[gv] (h) at (3.1,1.0) {$h$};
  \node[gv] (i) at (0.0,0.0) {$i$};
  \begin{scope}[on background layer]
  \draw[faint] (a) -- node[gw,above right]{1} (c);
  \draw[faint] (b) -- node[gw,above right]{2} (e);
  \draw[faint] (e) -- node[gw,left]{3} (h);
  \draw[nyss] (g) -- node[gw,below]{4} (h);
  \draw[cand] (b) -- node[gw,above left]{5} (d);
  \draw[ge] (d) -- node[gw,right]{6} (g);
  \draw[cand] (a) .. controls (-0.2,5.4) and (-0.4,1.3) .. node[gw,left,pos=0.5]{7} (g);
  \draw[ge] (g) -- node[gw,above left]{8} (i);
  \draw[faint] (a) -- node[gw,above left]{9} (b);
  \draw[faint] (e) -- node[gw,above right]{10} (f);
  \end{scope}
  \begin{scope}[on background layer]
  \fill[plum!15,rounded corners=10pt] (-0.55,-0.5) -- (1.55,0.6) -- (1.6,3.6) -- (0.55,3.75) -- (-0.55,1.5) -- cycle;
  \end{scope}
  \node[font=\scriptsize\bfseries,text=plum] at (0.55,4.1) {$S$};
\end{tikzpicture}
\par
\footnotesize Kanterna över snittet $S=\{d,g,i\}$: 4, 5, 7.
Den lättaste, \textbf{4}, ingår i \emph{varje} MST.
\end{column}
\begin{column}{0.49\textwidth}
\centering
{\small\bfseries\color{kthgreen}Cykelegenskapen}\par
\begin{tikzpicture}[scale=0.56,every node/.append style={font=\footnotesize}]
  \node[gv] (a) at (3.1,5.2) {$a$};
  \node[gv] (b) at (2.1,4.2) {$b$};
  \node[gv] (c) at (4.3,4.1) {$c$};
  \node[gv] (d) at (1.0,3.1) {$d$};
  \node[gv] (e) at (3.1,3.1) {$e$};
  \node[gv] (f) at (4.7,2.1) {$f$};
  \node[gv] (g) at (1.0,1.0) {$g$};
  \node[gv] (h) at (3.1,1.0) {$h$};
  \node[gv] (i) at (0.0,0.0) {$i$};
  \begin{scope}[on background layer]
  \draw[faint] (a) -- node[gw,above right]{1} (c);
  \draw[draw=kthgrass,line width=2pt] (b) -- node[gw,above right]{2} (e);
  \draw[draw=kthgrass,line width=2pt] (e) -- node[gw,left]{3} (h);
  \draw[draw=kthgrass,line width=2pt] (g) -- node[gw,below]{4} (h);
  \draw[draw=kthgrass,line width=2pt] (b) -- node[gw,above left]{5} (d);
  \draw[rejnu] (d) -- node[gw,right]{6} (g);
  \draw[faint] (a) .. controls (-0.2,5.4) and (-0.4,1.3) .. node[gw,left,pos=0.5]{7} (g);
  \draw[faint] (g) -- node[gw,above left]{8} (i);
  \draw[faint] (a) -- node[gw,above left]{9} (b);
  \draw[faint] (e) -- node[gw,above right]{10} (f);
  \end{scope}
\end{tikzpicture}
\par
\footnotesize På cykeln $b$--$d$--$g$--$h$--$e$ är \textbf{6} tyngst.
Den ingår i \emph{inget} MST.
\end{column}
\end{columns}
\vfill
\begin{kaninbox}[fontupper=\footnotesize]
Båda algoritmerna tar i varje steg en kant som är lättast över \emph{något} snitt som ingen vald
kant korsar. Därför är båda korrekta. Men snittet spelar olika roll: för Prim \emph{är} det
valregeln, för Kruskal finns det bara i beviset. Nästa bild.
\end{kaninbox}
\end{frame}}

{\kaninbg
\begin{frame}{\kh Snittet: valregel för Prim, bara bevis för Kruskal}
\begin{columns}[T]
\begin{column}{0.44\textwidth}
\centering
{\small\bfseries\color{kthblue}Kruskal efter två kanter}\par
\begin{tikzpicture}[scale=0.5,every node/.append style={font=\footnotesize}]
  \node[gv,comp1] (a) at (3.1,5.2) {$a$};
  \node[gv,comp2] (b) at (2.1,4.2) {$b$};
  \node[gv,comp1] (c) at (4.3,4.1) {$c$};
  \node[gv] (d) at (1.0,3.1) {$d$};
  \node[gv,comp2] (e) at (3.1,3.1) {$e$};
  \node[gv] (f) at (4.7,2.1) {$f$};
  \node[gv] (g) at (1.0,1.0) {$g$};
  \node[gv] (h) at (3.1,1.0) {$h$};
  \node[gv] (i) at (0.0,0.0) {$i$};
  \begin{scope}[on background layer]
  \fill[plum!28,rotate around={-47.7:(2.6,3.65)}] (2.6,3.65) ellipse (1.35 and 0.62);
  \draw[mst] (a) -- node[gw,above right]{1} (c);
  \draw[mst] (b) -- node[gw,above right]{2} (e);
  \draw[nyss] (e) -- node[gw,left]{3} (h);
  \draw[cand] (g) -- node[gw,below]{4} (h);
  \draw[cand] (b) -- node[gw,above left]{5} (d);
  \draw[cand] (d) -- node[gw,right]{6} (g);
  \draw[cand] (a) .. controls (-0.2,5.4) and (-0.4,1.3) .. node[gw,left,pos=0.5]{7} (g);
  \draw[cand] (g) -- node[gw,above left]{8} (i);
  \draw[cand] (a) -- node[gw,above left]{9} (b);
  \draw[cand] (e) -- node[gw,above right]{10} (f);
  \end{scope}
  \node[font=\small\bfseries,text=plum] at (2.2,2.95) {$S$};
\end{tikzpicture}
\par
{\scriptsize\textcolor{kthorange}{streckat}: kandidater\quad \textcolor{plum}{$S$}: bara i beviset}
\end{column}
\begin{column}{0.53\textwidth}
\footnotesize
\textbf{Valet är globalt.} Kruskal jämför \emph{alla} kanter mellan olika komponenter, 3 till 10
— också 4, 6, 7 och 8, som inte rör $b$ eller $e$. Den tar 3.\par\smallskip
\textbf{Snittet kommer i efterhand.} Beviset väljer $S$ = ena ändens komponent, här $\{b,e\}$.
Kanterna över $(S,V\setminus S)$ är 3, 5, 9 och 10. Alla går mellan olika komponenter, så alla var
kandidater. Alltså är 3 lättast även över snittet.\par\smallskip
Globalt lättast $\Rightarrow$ lättast över snittet. Inte tvärtom.
\end{column}
\end{columns}
\vfill
\footnotesize
\renewcommand{\arraystretch}{1.05}
\begin{tabularx}{\textwidth}{@{}p{0.12\textwidth}p{0.44\textwidth}X@{}}
\toprule
& \textbf{Vilken kant tas?} & \textbf{Snittet}\\
\midrule
Prim & lättast ut ur trädet & (trädet, resten): \emph{är} valregeln\\
Kruskal & lättast av alla kanter mellan olika komponenter & ena ändens komponent: bara i beviset\\
Borůvka & varje komponent tar sin lättaste kant ut & alla komponenters snitt, samtidigt\\
\bottomrule
\end{tabularx}
\end{frame}}

{\kaninbg
\begin{frame}{\kh Union--Find: hur Kruskal vet vad som hänger ihop}
\begin{columns}[T]
\begin{column}{0.52\textwidth}
\small
\begin{itemize}\setlength{\itemsep}{2pt}
  \item Varje komponent är ett träd av föräldrapekare; roten är komponentens namn.
  \item \textsc{Find}$(x)$: följ pekarna till roten.
  \item \textsc{Union}: häng den \emph{mindre} roten under den större.
  \item \textbf{Stigkomprimering:} låt alla på vägen peka direkt på roten.
  \item Båda knepen: amorterat $O(\alpha(n))$ per operation. $\alpha$ är i praktiken $\le4$.
\end{itemize}
\end{column}
\begin{column}{0.45\textwidth}
\centering
\begin{tikzpicture}[font=\small,
   uf/.style={circle,draw=kthnavy,minimum size=5mm,inner sep=0pt},
   p/.style={->,draw=kthblue,line width=0.8pt}]
  \node[font=\footnotesize\bfseries] at (0.7,3.3) {före \textsc{Find}$(z)$};
  \node[uf,fill=kthlight] (r) at (0.7,2.7) {$r$};
  \node[uf,fill=white] (x) at (0.7,1.9) {$x$};
  \node[uf,fill=white] (y) at (0.7,1.1) {$y$};
  \node[uf,fill=lampyellow!50] (z) at (0.7,0.3) {$z$};
  \draw[p] (x) -- (r); \draw[p] (y) -- (x); \draw[p] (z) -- (y);
  \draw[->,kthorange,line width=1pt] (1.35,1.5) -- (2.25,1.5);
  \node[font=\footnotesize\bfseries] at (3.4,3.3) {efter};
  \node[uf,fill=kthlight] (r2) at (3.4,2.7) {$r$};
  \node[uf,fill=white] (x2) at (2.6,1.6) {$x$};
  \node[uf,fill=white] (y2) at (3.4,1.6) {$y$};
  \node[uf,fill=lampyellow!50] (z2) at (4.2,1.6) {$z$};
  \draw[p] (x2) -- (r2); \draw[p] (y2) -- (r2); \draw[p] (z2) -- (r2);
\end{tikzpicture}
\end{column}
\end{columns}
\vfill
\begin{kaninbox}[fontupper=\footnotesize]
\textbf{Kruskals tid:} sortering $O(|E|\log|E|)$ + $|E|$ \textsc{Find} och $|V|-1$
\textsc{Union}. Sorteringen dominerar. Och $\log|E|\le\log|V|^2=2\log|V|$, så
$O(|E|\log|E|)=O(|E|\log|V|)$.
\end{kaninbox}
\end{frame}}

{\kaninbg
\begin{frame}{\kh Prims tid beror på prioritetskön}
\centering
\small
\begin{tabular}{@{}llll@{}}
\toprule
Kö & \textsc{ExtractMin} & \textsc{DecreaseKey} & Totalt\\
\midrule
osorterad array & $O(|V|)$ & $O(1)$ & $O(|V|^2)$\\
binär heap & $O(\log|V|)$ & $O(\log|V|)$ & $O(|E|\log|V|)$\\
Fibonacciheap & $O(\log|V|)$ am. & $O(1)$ am. & $O(|E|+|V|\log|V|)$\\
”lat” heap med kanter & $O(\log|E|)$ & — (dubbletter i stället) & $O(|E|\log|V|)$, se nästa bild\\
\bottomrule
\end{tabular}
\vfill
\begin{columns}[T]
\begin{column}{0.49\textwidth}
\begin{kaninbox}[fontupper=\footnotesize]
\textbf{Tät graf}, $|E|=\Theta(|V|^2)$: arrayen vinner. $O(|V|^2)$ är dessutom
\emph{linjärt i grannmatrisens storlek}.
\end{kaninbox}
\end{column}
\begin{column}{0.49\textwidth}
\begin{kaninbox}[fontupper=\footnotesize]
\textbf{Gles graf}, $|E|=O(|V|)$: heapen vinner, $O(|V|\log|V|)$.
Samma algoritm, olika datastruktur, olika ”bästa”.
\end{kaninbox}
\end{column}
\end{columns}
\varstrip{Separation}{algoritmen (Prim)}{datastrukturen och tätheten}{att komplexitet hör till algoritm \emph{plus} datastruktur}
\end{frame}}

{\kaninbg
\begin{frame}{\kh Lat heap: blir dubbletterna ett problem?}
\begin{columns}[T]
\begin{column}{0.47\textwidth}
{\small\bfseries\color{kthnavy}PrimLat$(G,w,r)$}
\footnotesize
\begin{algorithmic}[1]
\State $T\gets\{r\}$;\ $H\gets$ alla kanter $(r,x)$
\While{$H\neq\emptyset$}
  \State $(u,v)\gets\textsc{ExtractMin}(H)$
  \If{$v\in T$}
    \State släng den \Comment{inaktuell dubblett}
  \Else
    \State $T\gets T\cup\{v\}$; ta med $(u,v)$
    \ForAll{$(v,x)\in E$ med $x\notin T$}
      \State \textsc{Insert}$(H,(v,x))$
    \EndFor
  \EndIf
\EndWhile
\end{algorithmic}
\end{column}
\begin{column}{0.5\textwidth}
\footnotesize
\textbf{Ja,} i en fullständig graf ligger ett hörn i heapen en gång för varje granne som hann in
i trädet före det: upp till $|V|-1$ gånger.\par\medskip
\textbf{Men} varje kant läggs in högst en gång, när dess första ände kommer med i trädet. Heapen
har aldrig fler än $|E|$ element.\par\medskip
Alltså högst $|E|$ \textsc{Insert} och $|E|$ \textsc{ExtractMin}, var och en $O(\log|E|)$.
Och $\log|E|\le\log|V|^2=2\log|V|$.
\end{column}
\end{columns}
\vfill
\begin{columns}[T]
\begin{column}{0.47\textwidth}
\begin{ahabox}[fontupper=\footnotesize]
Totalt $O(|E|\log|V|)$, samma som med \textsc{DecreaseKey}. Dubbletterna kostar minne,
$O(|E|)$ i stället för $O(|V|)$, och en konstant faktor. Inte asymptotisk tid.
\end{ahabox}
\end{column}
\begin{column}{0.5\textwidth}
\begin{kaninbox}[fontupper=\footnotesize]
\textbf{Olyckliga vikter} behövs i varianten som bara lägger in när nyckeln blir bättre:
$w(i,j)=n(n-i)+j$ för $i<j$. Prim från 1 tar hörnen i ordning $1,2,\dots,n$, och varje nytt hörn
förbättrar alla återstående nycklar. Också då: $|E|$ element.
\end{kaninbox}
\end{column}
\end{columns}
\end{frame}}

{\kaninbg
\begin{frame}{\kh Borůvka (1926): alla komponenter väljer samtidigt}
\begin{columns}[T]
\begin{column}{0.49\textwidth}
\centering
{\small\bfseries\color{kthblue}Runda 1}\par
\input{gen/boruvka1}\par
\footnotesize Varje hörn tar sin billigaste kant:\\ 1, 2, 3, 4, 5, 8, 10. Två komponenter kvar.
\end{column}
\begin{column}{0.49\textwidth}
\centering
{\small\bfseries\color{kthblue}Runda 2}\par
\begin{tikzpicture}[scale=0.56,every node/.append style={font=\footnotesize}]
  \node[gv,comp2] (a) at (3.1,5.2) {$a$};
  \node[gv,comp2] (b) at (2.1,4.2) {$b$};
  \node[gv,comp2] (c) at (4.3,4.1) {$c$};
  \node[gv,comp2] (d) at (1.0,3.1) {$d$};
  \node[gv,comp2] (e) at (3.1,3.1) {$e$};
  \node[gv,comp2] (f) at (4.7,2.1) {$f$};
  \node[gv,comp2] (g) at (1.0,1.0) {$g$};
  \node[gv,comp2] (h) at (3.1,1.0) {$h$};
  \node[gv,comp2] (i) at (0.0,0.0) {$i$};
  \begin{scope}[on background layer]
  \draw[mst] (a) -- node[gw,above right]{1} (c);
  \draw[mst] (b) -- node[gw,above right]{2} (e);
  \draw[mst] (e) -- node[gw,left]{3} (h);
  \draw[mst] (g) -- node[gw,below]{4} (h);
  \draw[mst] (b) -- node[gw,above left]{5} (d);
  \draw[faint] (d) -- node[gw,right]{6} (g);
  \draw[nyss] (a) .. controls (-0.2,5.4) and (-0.4,1.3) .. node[gw,left,pos=0.5]{7} (g);
  \draw[mst] (g) -- node[gw,above left]{8} (i);
  \draw[cand] (a) -- node[gw,above left]{9} (b);
  \draw[mst] (e) -- node[gw,above right]{10} (f);
  \end{scope}
\end{tikzpicture}
\par
\footnotesize Billigaste kanten mellan dem: \textbf{7} (inte 9). Klart.
\end{column}
\end{columns}
\vfill
\begin{kaninbox}[fontupper=\footnotesize]
Antalet komponenter minst halveras per runda $\Rightarrow$ $O(\log|V|)$ rundor à $O(|E|)$:
$O(|E|\log|V|)$. Äldst av de tre, och lättast att parallellisera. (Kräver unika vikter
eller konsekvent tie-break — annars kan två komponenter välja varsin kant och bilda en cykel.)
\end{kaninbox}
\end{frame}}

% ---------------------------------------------------------------- variation
\begin{frame}{Vanligt fel: ”billigaste kanten från där jag står”}
\begin{columns}[T]
\begin{column}{0.42\textwidth}
\centering
\input{gen/girig}
\end{column}
\begin{column}{0.55\textwidth}
\begin{missbox}
”Prim = gå längs billigaste kanten från det senaste hörnet.”
\end{missbox}
\small
Från $a$: billigast är 1 till $c$. Kommer till $c$. Stannar i $c$.
\medskip

Det där är en \emph{promenad}, inte Prim. Prim tittar på alla kanter
ut ur \emph{hela trädet} $\{a,c\}$, inte bara ut ur $c$.
\medskip

Och även när promenaden inte kör fast blir den en stig — ett spännande träd behöver
kunna förgrena sig.
\end{column}
\end{columns}
\varstrip{Kontrast}{grafen, startpunkten}{”ut ur senaste hörnet” / ”ut ur trädet”}{snittet går runt hela trädet}
\end{frame}

\begin{frame}{Prim är inte Dijkstra}
\begin{columns}[T]
\begin{column}{0.33\textwidth}
\centering
{\small\bfseries\color{kthblue}MST (Prim), vikt 40}\par
\begin{tikzpicture}[scale=0.56,every node/.append style={font=\footnotesize}]
  \node[gv,intree] (a) at (3.1,5.2) {$a$};
  \node[gv] (b) at (2.1,4.2) {$b$};
  \node[gv] (c) at (4.3,4.1) {$c$};
  \node[gv] (d) at (1.0,3.1) {$d$};
  \node[gv] (e) at (3.1,3.1) {$e$};
  \node[gv] (f) at (4.7,2.1) {$f$};
  \node[gv] (g) at (1.0,1.0) {$g$};
  \node[gv] (h) at (3.1,1.0) {$h$};
  \node[gv] (i) at (0.0,0.0) {$i$};
  \begin{scope}[on background layer]
  \draw[mst] (a) -- node[gw,above right]{1} (c);
  \draw[mst] (b) -- node[gw,above right]{2} (e);
  \draw[mst] (e) -- node[gw,left]{3} (h);
  \draw[mst] (g) -- node[gw,below]{4} (h);
  \draw[mst] (b) -- node[gw,above left]{5} (d);
  \draw[faint] (d) -- node[gw,right]{6} (g);
  \draw[mst] (a) .. controls (-0.2,5.4) and (-0.4,1.3) .. node[gw,left,pos=0.5]{7} (g);
  \draw[mst] (g) -- node[gw,above left]{8} (i);
  \draw[faint] (a) -- node[gw,above left]{9} (b);
  \draw[mst] (e) -- node[gw,above right]{10} (f);
  \end{scope}
\end{tikzpicture}

\end{column}
\begin{column}{0.33\textwidth}
\centering
{\small\bfseries\color{plum}Kortaste-vägträd från $a$, vikt 47}\par
\begin{tikzpicture}[scale=0.56,every node/.append style={font=\footnotesize}]
  \node[gv,intree] (a) at (3.1,5.2) {$a$};
  \node[gv] (b) at (2.1,4.2) {$b$};
  \node[gv] (c) at (4.3,4.1) {$c$};
  \node[gv] (d) at (1.0,3.1) {$d$};
  \node[gv] (e) at (3.1,3.1) {$e$};
  \node[gv] (f) at (4.7,2.1) {$f$};
  \node[gv] (g) at (1.0,1.0) {$g$};
  \node[gv] (h) at (3.1,1.0) {$h$};
  \node[gv] (i) at (0.0,0.0) {$i$};
  \begin{scope}[on background layer]
  \draw[draw=plum,line width=2.6pt] (a) -- node[gw,above right]{1} (c);
  \draw[draw=plum,line width=2.6pt] (b) -- node[gw,above right]{2} (e);
  \draw[faint] (e) -- node[gw,left]{3} (h);
  \draw[draw=plum,line width=2.6pt] (g) -- node[gw,below]{4} (h);
  \draw[faint] (b) -- node[gw,above left]{5} (d);
  \draw[draw=plum,line width=2.6pt] (d) -- node[gw,right]{6} (g);
  \draw[draw=plum,line width=2.6pt] (a) .. controls (-0.2,5.4) and (-0.4,1.3) .. node[gw,left,pos=0.5]{7} (g);
  \draw[draw=plum,line width=2.6pt] (g) -- node[gw,above left]{8} (i);
  \draw[draw=plum,line width=2.6pt] (a) -- node[gw,above left]{9} (b);
  \draw[draw=plum,line width=2.6pt] (e) -- node[gw,above right]{10} (f);
  \end{scope}
\end{tikzpicture}
\par
{\scriptsize\color{plum}avstånd: $b$ 9, $c$ 1, $d$ 13, $e$ 11, $f$ 21, $g$ 7, $h$ 11, $i$ 15}
\end{column}
\begin{column}{0.3\textwidth}
\footnotesize
Samma skelett: kö, \textsc{ExtractMin}, \textsc{DecreaseKey}. En rad skiljer:\par\smallskip
Prim: $\mathit{key}[v]\gets w(u,v)$\par
Dijkstra: $d[v]\gets d[u]+w(u,v)$\par\smallskip
I MST är vägen $a\leadsto b$ lång: $7+4+3+2=16$. Dijkstra tar kanten 9 direkt.
MST minimerar \emph{summan över trädet}, inte avståndet från $a$.
\end{column}
\end{columns}
\varstrip{Kontrast}{graf, rot, köstruktur}{nyckeln: kantvikt eller avstånd}{vad som minimeras}
\end{frame}

\begin{frame}{Ändra vikterna: MST och kortaste-vägträd reagerar olika}
\scriptsize
\renewcommand{\arraystretch}{1.05}
\begin{tabularx}{\textwidth}{@{}p{0.24\textwidth}p{0.1\textwidth}X@{}}
\toprule
\textbf{Ändring} & \textcolor{kthblue}{\textbf{MST}} & \textcolor{plum}{\textbf{Kortaste-vägträd från $a$}} (förra bildens träd)\\
\midrule
$w\mapsto w^2$ & samma & \textbf{ändras}: $b$ nås nu via $e$ ($78<81$)\\
$w\mapsto w-5$ & samma & \textbf{odefinierat}: fram och tillbaka på en negativ kant\\
$w\mapsto w+c$ & samma & kan ändras: straffar stigar med många kanter\\
byt $5\leftrightarrow6$ på $b$--$d$, $d$--$g$ & \textbf{ändras} & samma (här)\\
\bottomrule
\end{tabularx}
\smallskip
\begin{columns}[T]
\begin{column}{0.32\textwidth}
\centering
{\scriptsize\bfseries\color{kthblue}MST med vikterna $w^2$}\par
\begin{tikzpicture}[scale=0.37,every node/.append style={font=\footnotesize}]
  \node[gv] (a) at (3.1,5.2) {$a$};
  \node[gv] (b) at (2.1,4.2) {$b$};
  \node[gv] (c) at (4.3,4.1) {$c$};
  \node[gv] (d) at (1.0,3.1) {$d$};
  \node[gv] (e) at (3.1,3.1) {$e$};
  \node[gv] (f) at (4.7,2.1) {$f$};
  \node[gv] (g) at (1.0,1.0) {$g$};
  \node[gv] (h) at (3.1,1.0) {$h$};
  \node[gv] (i) at (0.0,0.0) {$i$};
  \begin{scope}[on background layer]
  \draw[mst] (a) -- node[gw,above right]{1} (c);
  \draw[mst] (b) -- node[gw,above right]{4} (e);
  \draw[mst] (e) -- node[gw,left]{9} (h);
  \draw[mst] (g) -- node[gw,below]{16} (h);
  \draw[mst] (b) -- node[gw,above left]{25} (d);
  \draw[faint] (d) -- node[gw,right]{36} (g);
  \draw[mst] (a) .. controls (-0.2,5.4) and (-0.4,1.3) .. node[gw,left,pos=0.5]{49} (g);
  \draw[mst] (g) -- node[gw,above left]{64} (i);
  \draw[faint] (a) -- node[gw,above left]{81} (b);
  \draw[mst] (e) -- node[gw,above right]{100} (f);
  \end{scope}
\end{tikzpicture}
\par
{\scriptsize\color{kthblue}Samma kanter som förut.}
\end{column}
\begin{column}{0.32\textwidth}
\centering
{\scriptsize\bfseries\color{plum}Kortaste-vägträd från $a$, vikterna $w^2$}\par
\begin{tikzpicture}[scale=0.46,every node/.append style={font=\footnotesize}]
  \node[gv,intree] (a) at (3.1,5.2) {$a$};
  \node[gv] (b) at (2.1,4.2) {$b$};
  \node[gv] (c) at (4.3,4.1) {$c$};
  \node[gv] (d) at (1.0,3.1) {$d$};
  \node[gv] (e) at (3.1,3.1) {$e$};
  \node[gv] (f) at (4.7,2.1) {$f$};
  \node[gv] (g) at (1.0,1.0) {$g$};
  \node[gv] (h) at (3.1,1.0) {$h$};
  \node[gv] (i) at (0.0,0.0) {$i$};
  \begin{scope}[on background layer]
  \draw[draw=plum,line width=2.6pt] (a) -- node[gw,above right]{1} (c);
  \draw[draw=plum,line width=2.6pt] (b) -- node[gw,above right]{4} (e);
  \draw[draw=plum,line width=2.6pt] (e) -- node[gw,left]{9} (h);
  \draw[draw=plum,line width=2.6pt] (g) -- node[gw,below]{16} (h);
  \draw[faint] (b) -- node[gw,above left]{25} (d);
  \draw[draw=plum,line width=2.6pt] (d) -- node[gw,right]{36} (g);
  \draw[draw=plum,line width=2.6pt] (a) .. controls (-0.2,5.4) and (-0.4,1.3) .. node[gw,left,pos=0.5]{49} (g);
  \draw[draw=plum,line width=2.6pt] (g) -- node[gw,above left]{64} (i);
  \draw[faint] (a) -- node[gw,above left]{81} (b);
  \draw[draw=plum,line width=2.6pt] (e) -- node[gw,above right]{100} (f);
  \end{scope}
\end{tikzpicture}
\par
{\scriptsize\color{plum}\textbf{Inte} ett MST. Nytt: $b$ hänger på $e$.}
\end{column}
\begin{column}{0.32\textwidth}
\footnotesize
Strikt växande funktion av vikterna $\Rightarrow$ samma MST: Kruskal ser bara ordningen.\par\smallskip
Kortaste vägar jämför \emph{summor}: $7+4+3+2=16>9$, men $7^2+4^2+3^2+2^2=78<81=9^2$.
\end{column}
\end{columns}
\varstrip{Separation}{grafen, ändringen}{vilket träd}{MST följer ordningen, kortaste väg summorna}
\end{frame}

\begin{frame}{Är MST unikt?}
\begin{columns}[T]
\begin{column}{0.3\textwidth}
\centering
{\scriptsize\bfseries Om $d$--$g$ också väger 5 \dots}\par
\input{gen/tie1}
\end{column}
\begin{column}{0.3\textwidth}
\centering
{\scriptsize\bfseries \dots\ finns två MST, båda 40}\par
\input{gen/tie2}
\end{column}
\begin{column}{0.36\textwidth}
\footnotesize
\begin{itemize}\setlength{\itemsep}{3pt}
  \item Unika vikter $\Rightarrow$ exakt ett MST. (Bevis: cykelegenskapen.)
  \item Lika vikter: flera MST kan finnas, men alla har samma \emph{vikt}.
  \item Kruskal och Prim är fortfarande korrekta. Vilket träd de hittar beror på
        hur lika nycklar bryts.
\end{itemize}
\end{column}
\end{columns}
\vfill
\begin{missbox}
”Prim och Kruskal gav olika träd, alltså har jag räknat fel.” Bara om vikterna är unika.
\end{missbox}
\varstrip{Generalisering}{algoritmerna}{vikterna, från unika till lika}{”minimalt” betyder minsta vikt, inte ett visst träd}
\end{frame}

\begin{frame}{Mitt i körningen: går det att se vilken algoritm det är?}
\begin{columns}[T]
\begin{column}{0.33\textwidth}
\centering
{\small\bfseries\color{kthblue}Kruskal, 4 kanter}\par
\input{gen/skog}\par
\footnotesize Två träd med kanter, $\{a,c\}$ och $\{b,e,g,h\}$. Prim har alltid \emph{ett}:
inte Prim.
\end{column}
\begin{column}{0.33\textwidth}
\centering
{\small\bfseries\color{kthblue}Prim, 4 kanter}\par
\begin{tikzpicture}[scale=0.56,every node/.append style={font=\footnotesize}]
  \node[gv,intree] (a) at (3.1,5.2) {$a$};
  \node[gv] (b) at (2.1,4.2) {$b$};
  \node[gv,intree] (c) at (4.3,4.1) {$c$};
  \node[gv] (d) at (1.0,3.1) {$d$};
  \node[gv,intree] (e) at (3.1,3.1) {$e$};
  \node[gv] (f) at (4.7,2.1) {$f$};
  \node[gv,intree] (g) at (1.0,1.0) {$g$};
  \node[gv,intree] (h) at (3.1,1.0) {$h$};
  \node[gv] (i) at (0.0,0.0) {$i$};
  \begin{scope}[on background layer]
  \draw[mst] (a) -- node[gw,above right]{1} (c);
  \draw[ge] (b) -- node[gw,above right]{2} (e);
  \draw[mst] (e) -- node[gw,left]{3} (h);
  \draw[mst] (g) -- node[gw,below]{4} (h);
  \draw[ge] (b) -- node[gw,above left]{5} (d);
  \draw[ge] (d) -- node[gw,right]{6} (g);
  \draw[mst] (a) .. controls (-0.2,5.4) and (-0.4,1.3) .. node[gw,left,pos=0.5]{7} (g);
  \draw[ge] (g) -- node[gw,above left]{8} (i);
  \draw[ge] (a) -- node[gw,above left]{9} (b);
  \draw[ge] (e) -- node[gw,above right]{10} (f);
  \end{scope}
\end{tikzpicture}
\par
\footnotesize Ett träd med 7 men utan 2. Kruskal tar de lättaste kanterna som inte sluter en
cykel, och 2 sluter ingen: inte Kruskal.
\end{column}
\begin{column}{0.3\textwidth}
\footnotesize
\textbf{Gemensam invariant:} de valda kanterna är en delmängd av något MST.\par\smallskip
Här går det att se skillnad — av två \emph{olika} skäl. Flera träd avslöjar Kruskal. En
överhoppad billig kant avslöjar Prim.\par\smallskip
Att det \emph{alltid} går vore för starkt. Nästa bild.
\end{column}
\end{columns}
\varstrip{Kontrast}{grafen, antal kanter}{algoritmen}{flera träd, eller en överhoppad billig kant}
\end{frame}

\begin{frame}{Motexempel: Kruskal och Prim gör exakt samma sak}
\begin{columns}[T]
\begin{column}{0.4\textwidth}
\centering
\begin{tikzpicture}[scale=0.68,every node/.append style={font=\footnotesize}]
  \node[gv,intree] (r) at (0,1.0) {$r$};
  \node[gv] (p) at (1.4,2.1) {$p$};
  \node[gv] (q) at (2.9,1.3) {$q$};
  \node[gv] (s) at (1.4,0.0) {$s$};
  \node[gv] (t) at (2.9,-0.45) {$t$};
  \begin{scope}[on background layer]
  \draw[mst] (r) -- node[gw,above left] {1} (p);
  \draw[mst] (p) -- node[gw,above right] {2} (q);
  \draw[mst] (r) -- node[gw,below left] {3} (s);
  \draw[mst] (s) -- node[gw,below] {4} (t);
  \draw[ge] (r) -- node[gw,above] {5} (q);
  \draw[ge] (q) -- node[gw,right] {6} (t);
  \draw[ge] (s) -- node[gw,below right] {7} (q);
  \end{scope}
\end{tikzpicture}

\end{column}
\begin{column}{0.57\textwidth}
\footnotesize
\renewcommand{\arraystretch}{1.12}
\begin{tabular}{@{}lll@{}}
\toprule
Steg & Kruskal & Prim från $r$\\
\midrule
1 & $r$--$p$ (1) & $r$--$p$ (1)\\
2 & $p$--$q$ (2) & $p$--$q$ (2)\\
3 & $r$--$s$ (3) & $r$--$s$ (3)\\
4 & $s$--$t$ (4) & $s$--$t$ (4)\\
\bottomrule
\end{tabular}\par\smallskip
5, 6 och 7 sluter cykler. Samma tillstånd efter varje steg: Kruskals ”skog” är hela tiden
ett enda träd.
\end{column}
\end{columns}
\vfill
\begin{columns}[T]
\begin{column}{0.4\textwidth}
\begin{ahabox}[fontupper=\footnotesize]
\textbf{Flera träd} med kanter $\Rightarrow$ inte Prim. \textbf{Ett träd} $\Rightarrow$ kan vara
båda. (Prim från $t$ börjar med 4.)
\end{ahabox}
\end{column}
\begin{column}{0.57\textwidth}
\begin{missbox}[fontupper=\footnotesize]
”Skog = Kruskal, träd = Prim.” En skog kan vara ett enda träd — precis när Kruskals nästa kant
hela tiden sitter i trädet runt startpunkten.
\end{missbox}
\end{column}
\end{columns}
\varstrip{Separation}{algoritmerna}{grafen}{Prim ger alltid ett träd; Kruskal ibland också}
\end{frame}
